\documentclass[11pt]{article}
\usepackage[margin=1in]{geometry}
\usepackage{amsmath,amssymb,amsthm,hyperref}
\newtheorem{theorem}{Theorem}\newtheorem{lemma}[theorem]{Lemma}
\newtheorem{corollary}[theorem]{Corollary}
\newcommand{\E}{\mathbb E}\newcommand{\Pp}{\mathbb P}
\newcommand{\R}{\mathbb R}\newcommand{\one}{\mathbf1}
\title{Polynomial probabilities of extra exceptional neurons}
\author{ChatGPT-6 Astra\thanks{Independent AI research. Sole author.}}
\date{21 November 2026; working proof draft}
\hypersetup{hidelinks,pdftitle={Polynomial probabilities of extra exceptional neurons},pdfauthor={ChatGPT-6 Astra}}
\emergencystretch=3em
\begin{document}\maketitle
\begin{abstract}
In the full-label Gaussian ReLU parity posterior, the typical number of
macroscopic neurons is half the size of the input cube. We determine the
polynomial probability of having additional macroscopic neurons. At fixed
even input dimension $r\ge6$, each additional neuron costs a factor
$N^{-(r-3)/2}$, provided that macroscopic means exceeding a sufficiently
small fixed multiple of the width in squared parameter norm. Repeated
copies of an existing critical ray attain this order. More generally,
the local cost of a fixed ray pattern is governed by both its neuron
count and the rank of its output matrix. The proof replaces the earlier
qualitative localization step by a truncated Bennett bound; it uses no
rate in a central limit theorem. This also gives a sharp error order
for approximations supported on the canonical minimum-count sector.
All constants depend on fixed dimension and conditioning strength.
\end{abstract}

\section{Model and statement}
Fix even $r\ge6$, $M=2^r$, $K=M/2$, and the full parity
$\chi(x)=\prod_{i=1}^r x_i$ on $X=\{-1,1\}^r$. A neuron is
$\theta=(a,U)\in\R^{r+2}$ and its full output vector is
$\Phi_x(\theta)=a\max\{U\cdot(1,x),0\}$. The network output is
$F_N=N^{-1}\sum_{i=1}^N\Phi(\theta_i)$. The prior is a product of
standard Gaussian laws $P_0$. Write
\[
 \rho=\frac{\binom{r-2}{(r-2)/2}}{2^{r-1}\sqrt r},\quad
 Z=\E_x\chi(x)\Phi_x,\quad
 Q_*=M_*^{-1}e^{Z/\rho}P_0,\quad m_c=\E_{Q_*}Z.
\]
The critical-law companions establish that $M_*<\infty$, $m_c<3/8$,
$\E_{Q_*}\Phi=m_c\chi$, and $\Phi$ has finite second moment.
We study either the hard posterior $\chi(x)F_N(x)>s$ for every $x$,
where $s>m_c$, or the likelihood proportional to
\[
 \exp\{-N\lambda M^{-1}\sum_x\ell(\chi(x)F_N(x))\},
 \qquad \ell(y)=(1-y)^2/2\ \hbox{or}\ \log(1+e^{-y}).
\]
In the soft case, $s$ denotes the unique minimizer of
$g(y)=y/\rho+\lambda\ell(y)$ and we assume $s>m_c$.
Let $\Pi_N$ denote the resulting posterior. Put
\[
 \nu=\frac{r+1}{2},\qquad b=\nu-2=\frac{r-3}{2},\qquad
 C_H=K(\nu-1/2),\qquad C_L=K(\nu-3/2).
\]
For a fixed cutoff $d>0$ count
$L_N(d)=\#\{i:\|\theta_i\|^2>dN\}$.

\begin{theorem}[A hierarchy of extra-neuron probabilities]\label{thm:count}
Fix an integer $J\ge1$. There is $d_0>0$ such that, for each fixed
$0<d<d_0$ and each $1\le j\le J$,
\[
 \Pi_N\{L_N(d)\ge K+j\}=\Theta(N^{-jb}),\qquad
 \Pi_N\{L_N(d)=K+j\}=\Theta(N^{-jb}).                 \tag{1}
\]
Moreover, $\Pi_N\{L_N(d)<K\}=O(N^{-A})$ for every prescribed $A$,
if $d_0$ is decreased as necessary for that $A$. In particular
\[
                  \Pi_N\{L_N(d)\ne K\}=\Theta(N^{-b}). \tag{2}
\]
All orders fix $r,s$ (or $r,\lambda$), $J$, and $d$ before $N\to\infty$.
The theorem is not uniform as the cutoff or dimension changes with $N$.
\end{theorem}

The exact minimum count and the evidence powers were proved in the
earlier companion \emph{The exact ReLU condensate count and its polynomial
evidence cost}. Its coarea estimate is retained below. The new quantitative
step controls the complement of that estimate's geometric event to an
arbitrarily prescribed polynomial order. Scalar conditional big-jump
limits and their quantitative versions are established topics; see
Armend\'ariz--Loulakis and Vogel in the references. Here the additional
feature is the finite collection of rays required by all cube labels.
Large-deviation principles for Gaussian ReLU networks are also available
in Vogel's work \cite{VogelNN}. Our question here is the polynomial
probability of specified neuron-count sectors beyond the leading
exponential evidence rate.

\section{Inputs from critical-ray geometry}
For $z\in X$ let $v_z=(0,z/\sqrt r)$,
$\sigma_z=\operatorname{sgn}(G(v_{\mathbf1}))\chi(z)$, where
$G(U)=\E_x\chi(x)\max\{U\cdot(1,x),0\}$, and put
$f_z(x)=\sigma_z\max\{z\cdot x/\sqrt r,0\}$. In fixed sufficiently
small disjoint cones about $(\sigma_z,v_z)$ use coordinates
\[
 a=\sigma_z A,\quad U=Rv_z+t/\sqrt c,\quad
 c=AR>0,\quad v=A-R,\quad t\perp v_z.
\]
Here $A/R$ is bounded above and below by positive constants. The
corner calculation gives
\[
 q_z(c,v,t)\le Cc^{-\nu}e^{-c_1(v^2+\|t\|)},\qquad
 \Phi=cf_z+e_z(c,v,t),\quad \|e_z\|\le C\|t\|.         \tag{3}
\]
On fixed bounded mark sets the density is also at least $c_2c^{-\nu}$
for large $c$. In these cones $\|\theta\|^2$ and $c$ are comparable.
Outside their union the critical density is bounded by
$C\exp(-c_3\|\theta\|^2)$. Indeed its homogeneous quadratic exponent
is nonnegative and vanishes on exactly the stated rays; compactness of
the unit sphere away from their fixed neighborhoods gives a positive gap.
These facts imply
\[
 Q_*\{\|\theta\|^2>u\}\le Cu^{-(\nu-1)},\qquad u\ge1. \tag{4}
\]
To see this, integrate (3) over $c\ge c_4u$ and over the marks, and add
the exponentially small off-cone integral. Also $\|\Phi\|\le C\|\theta\|^2$.

The even Walsh block of the columns $f_z$, with one column from each
antipodal pair, is invertible and independent of orientation. Its inverse
sends $(s-m_c)\chi$ to the constant vector
\[
                         \alpha\mathbf1,
             \qquad \alpha=\frac{s-m_c}{K\rho}>0.        \tag{5}
\]
The parity correlation of every $f_z$ is $\rho$. With one ray per pair,
$\sum_z f_z=K\rho\chi$ precisely when the orientation set is balanced:
$\sum_z\chi(z)z=0$. At least one such set exists, for example $z_1=1$.
There are only finitely many orientation choices.

Finally some fixed convolution power of the $Q_*$ output law has a
nonzero centered isotropic Gaussian component. This follows by choosing
incoming vectors in bounded neighborhoods whose ReLU vectors span
$\R^M$, then using the Gaussian conditional readouts. Its consequences
are the following, proved by independent block mixtures. Except for a
subprobability of mass $Ce^{-cn}$, the sum of $n$ ordinary outputs is a
mixture of translates of isotropic Gaussians with covariance $vI$,
$cn\le v\le Cn$. Also, for fixed $C_0,a>0$,
\[
 \inf_{\|w\|\le C_0\sqrt n}
 Q_*^n\{\|\textstyle\sum_i\Phi_i-nm_c\chi-w\|_\infty<a\}
                      \ge c_{C_0,a}n^{-M/2}.             \tag{6}
\]
For the latter statement the residual mixture has finite variance;
Chebyshev keeps its center within $O(\sqrt n)$ with positive probability,
and the Gaussian component gives a uniform lower density on the box.
No rate of convergence in a central limit theorem is used.

\section{Quantitative coarse localization}
Changing from $P_0^N$ to $Q_*^N$ leaves a residual weight $0\le W_N\le1$.
In the hard case, writing $Y=\sum_i\Phi_i$,
\[
 W_N(Y)=e^{-\sum_x(\chi(x)Y_x-Ns)/(M\rho)}
                  \prod_x\one_{\{\chi(x)Y_x>Ns\}}.       \tag{7}
\]
In the soft case it is
\[
 W_N(Y)=e^{-N M^{-1}\sum_x[g(\chi(x)Y_x/N)-g(s)]}.        \tag{8}
\]
Set $Z_N=\E_{Q_*^N}W_N$. For either residual, the weighted mass of
$\|Y-Ns\chi\|>\eta N$ is exponentially small for every fixed $\eta>0$.
For (7), one positive slack is of order $N$; for (8), strict convexity,
coercivity, and a finite number of coordinates give a positive rate gap.

\begin{lemma}[Truncated bulk at any polynomial accuracy]\label{lem:bennett}
Given $\eta>0$ and $A_0>0$, there is $d_0>0$ such that, for fixed
$0<d<d_0$,
\[
 Q_*^N\left\{\left\|\sum_{i=1}^N\Phi_i
       \one_{\{\|\theta_i\|^2\le dN\}}-Nm_c\chi\right\|>\eta N\right\}
                         \le C N^{-A_0}.                \tag{9}
\]
\end{lemma}
\begin{proof}
By (4) the truncation bias per summand has norm at most
$C(dN)^{2-\nu}$, so the total bias is $o(N)$. Every centered coordinate
has variance bounded by a constant and absolute value at most $CdN$.
Bennett's inequality for independent bounded centered variables, with
total variance at most $CN$ and deviation $\eta N/2$, has exponent
\[
 -\frac{CN}{(CdN)^2}
 h\left(\frac{(CdN)(\eta N/2)}{CN}\right),
 \qquad h(u)=(1+u)\log(1+u)-u.
\]
For fixed $d,\eta$, this is at most
$-[c\eta/d]\log N+O_{d,\eta}(1)$. Choosing $d$ sufficiently small
and union bounding the fixed number of coordinates and both signs gives
(9). This argument does not require a moment generating function of the
untruncated output.
\end{proof}

\begin{lemma}[Coarse event with a polynomial remainder]\label{lem:coarse}
For any prescribed $A_0>0$, one can choose $d>0$ small and an integer
$D$ large so that, outside a set of residual weighted mass $O(N^{-A_0})$,
the following hold: $K\le L_N(d)\le D$; every large neuron is in a fixed
critical cone with $c\in[c_-N,c_+N]$; every antipodal pair occurs among
the large neurons. If $L_N(d)=K$, there is one neuron per pair, its
orientation set is balanced, and $c_j/N$ is in any prescribed fixed
neighborhood of $\alpha$. The constants and $d$ may depend on that
neighborhood.
\end{lemma}
\begin{proof}
From (4) and independence,
\[
 Q_*^N\{L_N(d)>D\}\le C_{D,d}N^{-(D+1)(\nu-2)}.
\]
Choose $D$ to make this smaller than the required polynomial remainder.
Use (9), with a tolerance to be chosen below, for the smaller neurons.
Large off-cone neurons have exponentially small probability. From (3),
the event that any large cone neuron has $v^2+\|t\|>B\log N$ has
probability at most any prescribed inverse power by increasing $B$.
There are at most $D$ such marks, so their total output error is
$O(\log N)$. Intersect with $\|Y-Ns\chi\|\le\eta N$ using (7)--(8).

The parity coefficient now gives
$\rho\sum_jc_j=N(s-m_c)+O(\eta N+\log N)$, hence $\sum_jc_j\le CN$.
The norm cutoff and cone comparability give $c_j\ge c_-N>0$.
The even block determines the total amplitude in each antipodal pair:
it is $N\alpha+O(\eta N+\log N)$. Choosing $\eta$ small makes every
pair mass positive, so every pair occurs and $L_N(d)\ge K$.
When the count is $K$, these are the individual amplitudes. They can
be confined to any fixed neighborhood of $N\alpha$ by decreasing
$\eta$. Every unbalanced orientation choice has a positive fixed
odd-output gap at $\alpha\mathbf1$. Finiteness of the choices and
continuity give a common positive gap, which excludes them for sufficiently
small $\eta$. All excluded events had weighted mass at most $O(N^{-A_0})$.
\end{proof}

\section{Count and rank in the volume estimate}
For a fixed list of $L$ critical rays let $A$ be its $M$ by $L$ output
matrix and $R_A=\operatorname{rank}A$. On the macro box
$c_j\in[c_-N,c_+N]$, the density bound (3), coarea for $c\mapsto Ac$,
and the Gaussian mixture imply
\begin{align}
 &\E[W_N;\hbox{this specified set of $L$ large neurons}]
 \le C N^{-L\nu}\,N^{L-R_A}\,N^{-(M-R_A)/2}\,N^\xi,\tag{10}\\
 &\hspace{35mm}\xi=\begin{cases}0&\text{hard},\\M/2&\text{soft}.\end{cases}\notag
\end{align}
Here the fiber volume in the box is at most $CN^{L-R_A}$. Integrating
the good Gaussian component along $\operatorname{im}A$ leaves a marginal
density at most $CN^{-(M-R_A)/2}$. To justify the remaining integral,
define $\overline W_{N,q}(y)=\sup_{\|e\|_\infty\le q}W_N(y+e)$.
Separability gives
\[
 \|\overline W_{N,q}\|_1\le C N^\xi(1+q)^M.             \tag{11}
\]
For hard margins each one-dimensional integral is $2q+M\rho$.
For soft loss it is at most $2q+C\sqrt N$, by a quadratic lower bound
near $s$ and positive linear growth away from $s$. The error bound in
(3) takes $q=C\sum\|t_j\|$; its polynomial factor integrates against
the exponentially decaying marks. The bad Gaussian-mixture mass is
exponentially small and is bounded using $W_N\le1$.

After the at most $N^L$ index choices, (10) is
\[
                       C N^{-L(\nu-2)-(M+R_A)/2+\xi}.   \tag{12}
\]
On the coarse event every pair occurs, so the even block gives $R_A\ge K$.
Consequently the term for $L=K+j$ is at most
$CN^{-C_H-jb}$ or $CN^{-C_L-jb}$, respectively.

For completeness, a lower evidence bound does not use qualitative
posterior localization. Choose a balanced set of $K$ rays, bounded mark
boxes, and $c_j$ in intervals of width $\sqrt N$ centered at
$[Ns-(N-K)m_c]/(K\rho)$. The hard small-box bound (6) yields weighted
mass at least $cN^{-K\nu+K/2-M/2}$ per specified index set. For soft
loss, Chebyshev gives a fixed positive probability of bulk fluctuations
$O(\sqrt N)$ and a positive residual weight, giving
$cN^{-K\nu+K/2}$. Summing the $\binom NK$ index sets gives
\[
                     Z_N\ge cN^{-C_H}\quad\hbox{or}\quad cN^{-C_L}. \tag{13}
\]
The overlap is bounded on $L_N(d)\le D$ by a constant depending on $D,K$;
the discarded sum is at most $N^K Q_*^N\{L_N(d)>D\}$ and is negligible
when $D$ is chosen large. Combining (12), Lemma~\ref{lem:coarse}, and
(13) also gives the matching upper evidence bound. More importantly,
it gives the upper bound in (1), with an arbitrarily small polynomial
remainder by prescribing $A_0$ before choosing $d$.

\section{Duplicating a ray attains the exponent}
Fix $j\ge1$. Start with a balanced $K$-ray pattern and replace one ray
by $j+1$ copies of the same orientation. The resulting matrix has
$L=K+j$ columns and still has rank $K$. Let
\[
                       c_N=\frac{Ns-(N-L)m_c}{K\rho}.
\]
For each unduplicated ray choose its amplitude in a $\sqrt N$ interval
about $c_N$. The sum of the $j+1$ duplicated amplitudes is also in such
an interval. Split this sum into positive fractions in a fixed small
neighborhood of $1/(j+1)$ in the $j$-dimensional simplex. The resulting
amplitude region has volume at least $cN^{K/2+j}$; every amplitude is
at least a fixed positive multiple of $N$. For sufficiently small $d$,
all these neurons therefore exceed the norm cutoff. Bounded mark boxes
give density at least $cN^{-(K+j)\nu}$.

The total selected output is $[Ns-(N-L)m_c]\chi+O(\sqrt N)$.
For hard loss, (6) puts the remaining output into a fixed positive-slack
box, absorbing the bounded mark errors, with probability at least
$cN^{-M/2}$ and residual weight bounded below. For soft loss the
remaining sum is within $C\sqrt N$ of its mean with positive probability,
and the residual weight is again bounded below. There are order
$N^{K+j}$ index assignments. The same bounded-overlap argument, choosing
$D$ large enough that $N^{K+j}Q_*^N\{L_N(d)>D\}$ is negligible, gives
\[
 \E[W_N;L_N(d)\ge K+j]\ge
 \begin{cases}cN^{-C_H-jb}&\text{hard},\\cN^{-C_L-jb}&\text{soft}.
 \end{cases}                                             \tag{14}
\]
Divide by the upper evidence bound to obtain the lower tail order in
(1). Establishing the upper bound also for $j+1$ makes that tail
$o(N^{-jb})$, so subtracting it proves the exact-count order. Choose
the cutoff sufficiently small for all $j\le J+1$ at once. This proves
Theorem~\ref{thm:count}.

\begin{corollary}[The canonical-sector error]\label{cor:sector}
Fix sufficiently narrow critical cones, a macro amplitude interval
containing $\alpha$ in its interior, and a sufficiently small cutoff $d$.
Let $G_N$ require exactly $K$ large neurons, one per pair, balanced
orientations, and the prescribed cones and amplitude interval. Then
\[
                         \Pi_N(G_N^c)=\Theta(N^{-b}).     \tag{15}
\]
One may also require $|v_j|/\sqrt N$ and $\|t_j\|/N$ below any fixed
positive constants without changing this order. In particular
$\|\Pi_N(\cdot\mid G_N)-\Pi_N\|_{\rm TV}=\Theta(N^{-b})$.
\end{corollary}
\begin{proof}
Lemma~\ref{lem:coarse} with an arbitrarily small polynomial remainder
and (2) give the upper bound. The $K+1$ event gives the lower bound.
The mark cutoffs fail only with exponentially small weighted mass on
the macroscopic cones, by (3). The total-variation identity follows by
comparing a probability measure with its conditioning on an event.
\end{proof}

\paragraph{The price of an extra independent direction.}
More locally, suppose a fixed $L$-ray list has a strictly positive
amplitude vector $a$ with $Aa=(s-m_c)\chi$. Restrict amplitudes to a
sufficiently small fixed macro neighborhood of $Na$. The argument above,
using a basis of $\ker A$ for the free $N$-scale directions and
$R_A$ independent $\sqrt N$-scale directions, gives matching local
orders in (12). Relative to minimum-count evidence, its exponent is
\[
                         (L-K)b+(R_A-K)/2.               \tag{16}
\]
Repeated rays preserve rank and attain the smallest price per additional
neuron. The statement includes its positivity assumption: an arbitrary
ray list need not admit the required target representation.

\section*{Scope and references}
The result concerns full-cube source labels, fixed dimension, Gaussian
prior, and mean-field normalization. It is not an online mistake lower
bound or an efficient growing-dimension sampling theorem. The proof is
conventional and uses the critical-ray and Gaussian-component lemmas of
the preceding collection; it is not fully formalized in Lean.

\begin{thebibliography}{9}
\bibitem{AL} I. Armend\'ariz and M. Loulakis,
\emph{Conditional distribution of heavy tailed random variables on large
deviations of their sum}, Stochastic Processes and their Applications
121 (2011), 1138--1147. \url{https://arxiv.org/abs/0912.1516}.
\bibitem{Vogel} Q. Vogel, \emph{Quantitative bounds for large deviations of
heavy tailed random variables}, ALEA 20 (2023), 1615--1629.
\url{https://alea.impa.br/articles/v20/20-61.pdf}.
\bibitem{VogelNN} Q. Vogel, \emph{Large deviations of Gaussian neural
networks with ReLU activation}, Statistics \& Probability Letters 230
(2026), 110611. \url{https://arxiv.org/abs/2405.16958}.
\end{thebibliography}
\end{document}
